% Copyright (c) 2026 Mika Nystrom
% SPDX-License-Identifier: Apache-2.0

\documentclass[11pt]{article}

% Page layout
\usepackage[letterpaper, margin=1in]{geometry}
\setlength{\headheight}{14pt}
\usepackage{parskip}
\setlength{\parindent}{0pt}

% Typography
\usepackage{fancyvrb}
\usepackage{booktabs}
\usepackage{amsmath}
\usepackage{amssymb}

% Code listings
\usepackage{listings}
\lstset{
    basicstyle=\ttfamily\small,
    columns=fullflexible,
    keepspaces=true,
    showstringspaces=false,
    breaklines=true,
    frame=none,
    xleftmargin=2em,
}

% Hyperlinks
\usepackage[colorlinks=true, linkcolor=blue, urlcolor=blue, citecolor=blue]{hyperref}

% Headers and footers
\usepackage{fancyhdr}
\pagestyle{fancy}
\fancyhf{}
\fancyhead[L]{\textit{BDD Memory Optimization for Symbolic Deadlock Checking}}
\fancyhead[R]{\thepage}
\renewcommand{\headrulewidth}{0.4pt}

\title{\textbf{BDD Memory Optimization for Symbolic Deadlock Checking} \\[0.5em]
\large Approaches, Trade-offs, and Implementation Plan}
\author{Mika Nystr\"om \and Claude (Anthropic)}
\date{DRAFT\\
\today}

\begin{document}

\maketitle

\tableofcontents
\newpage

%======================================================================
\section{Introduction}
%======================================================================

The BDD-based symbolic deadlock checker for the Fulcrum CSP toolchain
(\texttt{symbolic.scm}) works correctly but uses excessive memory on
large problems.  For $N = 8$ dining philosophers (16 process instances,
32 BDD variables), peak resident memory is approximately 2.7\,GB and
execution takes 35\,seconds on Apple~M2.  The reachable state set
converges in 8~iterations with a compact BDD of only 1,564 nodes, yet
the process of \emph{computing} reachability consumes three orders of
magnitude more memory than the final result requires.

This report analyzes the root cause of the memory blowup, surveys six
approaches for reducing it, and recommends a phased implementation plan.

%======================================================================
\section{Root Cause Analysis}
%======================================================================

The current implementation builds a monolithic transition relation
$T(V, V')$ as a single BDD and computes the image at each iteration via:
%
\begin{equation}
\mathit{Image}(R, T) =
  \bigl(\exists V.\; R(V) \wedge T(V, V')\bigr)[V'/V]
\label{eq:image}
\end{equation}
%
The intermediate BDD $R(V) \wedge T(V, V')$ is the product of the
reached set (1,564 nodes at convergence) and the transition relation
(2,307 nodes).  Despite both operands being modest, their conjunction
explodes because it must represent, for every reachable state, the
full set of successor states---a much richer structure than either
operand alone.

This intermediate is live throughout the subsequent existential
quantification (which eliminates current-state variables one at a time
via 16 applications of $\exists v.\; f = f|_{v=1} \vee f|_{v=0}$)
and the variable renaming (16 substitutions).  The peak occurs at
the conjunction, before any variables have been quantified out.

A mark-and-sweep garbage collector was added to the BDD library
(\texttt{ClearCaches}, \texttt{MarkRoot}, \texttt{CollectGarbage})
and benchmarked.  Running GC between iterations had no effect on peak
memory (2,749\,MB with GC vs.\ 2,702\,MB without) and added 38\%
time overhead (47.8\,s vs.\ 34.7\,s).  The reason is clear:
the GC runs between iterations, but the peak occurs \emph{within}
an iteration, during the conjunction.

The GC hooks are retained in the BDD library for future use; the
question is what algorithmic changes can reduce the within-iteration peak.

%======================================================================
\section{Approach~1: Frontier-Based Image Computation}
%======================================================================

%----------------------------------------------------------------------
\subsection{Idea}
%----------------------------------------------------------------------

Standard BFS reachability computes $\mathit{Image}(R_k, T)$ where
$R_k$ is the cumulative set of all states reached so far.  As $R_k$
grows, the conjunction $R_k \wedge T$ grows with it, even though most
of the states in $R_k$ have already been expanded in previous iterations
and contribute nothing new.

Frontier-based reachability maintains a separate \emph{frontier} set
$F_k$ containing only the states discovered in iteration~$k$:
%
\begin{align}
F_0 &= \mathit{Init} \\
I_k &= \mathit{Image}(F_k, T) \\
F_{k+1} &= I_k \setminus R_k \\
R_{k+1} &= R_k \cup F_{k+1}
\end{align}
%
The fixed point is reached when $F_{k+1} = \emptyset$.

%----------------------------------------------------------------------
\subsection{Expected Benefit}
%----------------------------------------------------------------------

The frontier $F_k$ is typically much smaller than $R_k$, especially
in later iterations where most reachable states have already been
discovered.  The conjunction $F_k \wedge T$ therefore involves a
smaller left operand, reducing the peak intermediate BDD size.

For the $N = 8$ dining philosophers, $R$ peaks at 1,564 BDD nodes.
The frontier at each iteration is a subset of the newly discovered
states and is expected to be 2--5$\times$ smaller than $R$ in later
iterations.  The conjunction intermediate should shrink proportionally.

%----------------------------------------------------------------------
\subsection{Caveat: The Cost of Set Difference}
%----------------------------------------------------------------------

Computing $F_{k+1} = I_k \setminus R_k$ requires
$I_k \wedge \neg R_k$.  In the current BDD library, $\neg R_k$
creates a full-size BDD (the library has no complement edges).
This partially offsets the savings from using a smaller frontier.

However, the \texttt{And} operation short-circuits on many paths
(returning \texttt{false} whenever one operand is \texttt{false}),
so the actual intermediate created by $I_k \wedge \neg R_k$ is
typically smaller than $|I_k| + |\neg R_k|$ would suggest.

Furthermore, after computing $F_{k+1}$ and $R_{k+1}$, the GC
can reclaim the intermediate $\neg R_k$ and $I_k$, bounding the
accumulation of dead nodes across iterations.

%----------------------------------------------------------------------
\subsection{Implementation}
%----------------------------------------------------------------------

The change is approximately 20 lines in the reachability loop
of \texttt{symbolic.scm}, with no BDD library modifications.
The existing \texttt{BDD.CollectGarbage} can be called between
iterations to reclaim dead frontier and image BDDs.

%======================================================================
\section{Approach~2: Partitioned Transition Relations}
%======================================================================

%----------------------------------------------------------------------
\subsection{Idea}
%----------------------------------------------------------------------

Instead of forming a single monolithic $T = T_1 \vee T_2 \vee \cdots
\vee T_m$ (one disjunct per transition), keep the partitions separate
and compute the image incrementally:
%
\begin{equation}
\mathit{Image}(R, T) =
  \bigcup_{i=1}^{m} \mathit{Image}(R, T_i)
\label{eq:part-image}
\end{equation}
%
This is \emph{disjunctive partitioning}, which maps naturally to CSP
systems where each partition is either a single process's tau
transition or a channel synchronization between two processes.

%----------------------------------------------------------------------
\subsection{Early Quantification}
%----------------------------------------------------------------------

The critical optimization is \emph{early quantification}.  Each
partition $T_i$ mentions only a small subset of the state variables
(those belonging to the 1--2 participating process instances).
Variables that appear in $R$ but not in $T_i$ can be existentially
quantified out of $R$ \emph{before} forming the conjunction:
%
\begin{equation}
\mathit{Image}(R, T_i) =
  \bigl(\exists V_{\mathrm{early}}.\; R\bigr) \wedge T_i
  \quad\text{(then quantify remaining $V$ and rename)}
\end{equation}
%
where $V_{\mathrm{early}}$ are the current-state variables not
mentioned by $T_i$.

For a tau transition of process instance~$j$ in an $N$-instance
system with $b$~bits per instance, $T_i$ mentions only $2b$ variables
(current and next state of instance~$j$), plus frame conditions
that equate current and next state for all other instances.
After early quantification, $R_{\mathrm{projected}}$ has only $b$
BDD variables instead of $Nb$, making it dramatically smaller.

%----------------------------------------------------------------------
\subsection{Frame Condition Simplification}
%----------------------------------------------------------------------

The frame condition for non-participating instances
($\mathit{next}_k = \mathit{curr}_k$ for instance $k \neq j$)
becomes an identity under existential quantification of $\mathit{curr}_k$:
$\exists \mathit{curr}_k.\; (\mathit{curr}_k \Leftrightarrow
\mathit{next}_k) = \mathit{true}$.  This means the frame conditions
for non-participating instances can be omitted entirely from the
partitioned image computation, further simplifying each partition.

For channel synchronizations involving two instances $j_1$ and $j_2$,
the frame conditions for all other instances can similarly be dropped,
and early quantification reduces $R$ to the $2b$ variables of
$j_1$ and $j_2$.

%----------------------------------------------------------------------
\subsection{Expected Benefit}
%----------------------------------------------------------------------

For $N = 8$ dining philosophers with $b = 2$ bits per instance:

\begin{itemize}
\item The monolithic $T$ has 2,307 BDD nodes over 32 variables.
\item Each tau partition has $\sim$10--30 nodes over 4 variables.
\item Each channel sync partition has $\sim$20--50 nodes over 8 variables.
\item The projected $R$ for a single-instance tau partition has
  2 BDD variables, giving $|R_{\mathrm{projected}}| \leq 5$ nodes.
\item The conjunction $R_{\mathrm{projected}} \wedge T_i$ is tiny
  ($\sim$50 nodes), compared to the monolithic
  $R \wedge T$ which can reach hundreds of thousands of nodes.
\end{itemize}

The savings are multiplicative: each partition's image computation
is cheap, and there are $O(N + C)$ partitions (where $C$ is the
number of channels).  The total work per iteration is
$O((N + C) \cdot \text{small})$ instead of $O(\text{huge})$.

\textbf{Expected reduction: 5--20$\times$ in peak memory.}

%----------------------------------------------------------------------
\subsection{Implementation}
%----------------------------------------------------------------------

Changes are entirely in \texttt{symbolic.scm} (2--4 days of work):

\begin{enumerate}
\item Modify \texttt{build-transition-bdd} to return a list of
  partition BDDs (one per tau edge, one per channel sync) instead
  of their disjunction.  Each partition carries metadata identifying
  which instance indices it involves.
\item Rewrite \texttt{symbolic-image} to iterate over partitions,
  performing early quantification for each.
\item Simplify frame conditions: omit non-participating instance
  frames from each partition.
\end{enumerate}

No BDD library changes are required.  Combines naturally with
frontier-based computation (Approach~1) for multiplicative savings.

%======================================================================
\section{Approach~3: Bounded Model Checking with SAT}
%======================================================================

%----------------------------------------------------------------------
\subsection{Idea}
%----------------------------------------------------------------------

Bounded model checking (BMC)~\cite{biere1999} replaces BDD-based
fixpoint computation with propositional satisfiability (SAT).  The
transition relation is \emph{unrolled} for $k$ steps:
%
\begin{equation}
\mathit{Init}(V_0) \wedge
\bigwedge_{i=0}^{k-1} T(V_i, V_{i+1}) \wedge
\mathit{Bad}(V_k)
\label{eq:bmc}
\end{equation}
%
where $\mathit{Bad}(V_k) = \neg\exists V'.\; T(V_k, V')$ is the
deadlock condition.  If the formula is satisfiable, the satisfying
assignment directly encodes a counterexample trace of length~$k$.

%----------------------------------------------------------------------
\subsection{Advantages}
%----------------------------------------------------------------------

\begin{itemize}
\item \textbf{Memory:}  SAT solvers use memory linear in the clause
  database size.  For $N = 8$ with $b = 2$ bits and $k = 8$ (the
  reachability diameter), the SAT instance has $\sim$20 $\times$
  8 = 160 propositional variables and a moderate number of clauses.
  A modern CDCL solver (MiniSat, CaDiCaL) handles this in
  milliseconds with negligible memory.
\item \textbf{Counterexamples:}  BMC naturally produces minimum-length
  counterexamples (by trying $k = 0, 1, 2, \ldots$).
\item \textbf{No BDD library needed:}  The SAT solver is a separate
  engine; the BDD library is not involved.
\end{itemize}

%----------------------------------------------------------------------
\subsection{Limitations}
%----------------------------------------------------------------------

\begin{itemize}
\item \textbf{Deadlock detection only:}  BMC finds deadlocks reachable
  within $k$ steps.  To \emph{prove} deadlock freedom, one must either
  know the system's diameter (which requires computing full
  reachability) or use $k$-induction or property-directed
  reachability (IC3/PDR), both of which add significant complexity.
\item \textbf{Diameter unknown:}  Without a diameter bound, BMC must
  increase $k$ until either a deadlock is found or a resource limit
  is hit.  For deadlock-free systems, BMC alone cannot terminate with
  a proof.
\end{itemize}

BMC therefore complements rather than replaces the BDD-based checker:
use BMC for fast deadlock \emph{detection}, and BDDs for deadlock
\emph{freedom proofs}.

%----------------------------------------------------------------------
\subsection{Implementation}
%----------------------------------------------------------------------

The implementation requires:
\begin{enumerate}
\item A SAT solver.  Options include linking a C solver (MiniSat is
  $\sim$2000 lines) or implementing a simple DPLL/CDCL solver in
  Scheme for small instances.
\item A CNF encoder that translates each LTS transition into clauses
  over the unrolled state variables.
\item A driver loop that increases $k$ until a satisfying assignment
  is found.
\end{enumerate}

Estimated effort: 3--5 days.  No changes to the BDD library.

%======================================================================
\section{Approach~4: BDD Variable Reordering}
%======================================================================

%----------------------------------------------------------------------
\subsection{Static Ordering}
%----------------------------------------------------------------------

BDD size is exponentially sensitive to variable ordering.  The current
implementation interleaves current-state and next-state variables
\emph{within} each process instance:
%
\[
  c_{0,0},\; n_{0,0},\; c_{0,1},\; n_{0,1},\;\;
  c_{1,0},\; n_{1,0},\; c_{1,1},\; n_{1,1},\;\;
  \ldots
\]
%
This is already a good ordering: interleaving keeps frame conditions
($c_{i,k} \Leftrightarrow n_{i,k}$) compact, and sequential
instantiation places communicating processes (fork~$i$, philosopher~$i$)
at adjacent positions.

Static heuristics such as FORCE~\cite{force} (iteratively moving each
variable toward the center of gravity of its neighbors in the
constraint graph) can refine the ordering, but for the regular
dining-philosophers topology the improvement is expected to be
10--30\%.

%----------------------------------------------------------------------
\subsection{Dynamic Sifting}
%----------------------------------------------------------------------

Dynamic sifting~\cite{rudell1993} moves each variable to every
possible position, measuring total BDD size, and places it at the
minimum.  This is the most effective reordering technique but requires
deep integration with the BDD library: swapping adjacent variables
means rebuilding unique tables, rewriting node children, and
invalidating caches.

The current BDD library uses separate \texttt{Root} objects per
variable, with each \texttt{Root} holding its own unique table and
operation caches.  Implementing sifting would require:
\begin{enumerate}
\item A procedure to swap two adjacent variables in the ordering.
\item Rebuilding the unique tables of both variables (every node
  belonging to one variable whose children belong to the other must
  be reconstructed).
\item Invalidating all operation caches (the existing
  \texttt{ClearCaches} suffices).
\end{enumerate}

Estimated effort: 1--2 weeks.  Improvement: 10--50\%, highly
problem-dependent.

%----------------------------------------------------------------------
\subsection{Assessment}
%----------------------------------------------------------------------

Variable reordering is a well-understood technique with modest
returns for regular systems like dining philosophers.  It becomes
important for irregular systems (pipelines, buses, complex protocols).
Static heuristics are cheap to implement (1--2 days, Scheme only);
dynamic sifting is a significant BDD library modification for
moderate gains.

%======================================================================
\section{Approach~5: Complement Edges}
%======================================================================

In a standard BDD library, $\neg f$ creates a new BDD by
recursively complementing every terminal, doubling the node count
in the worst case.  With complement edges (also called negation arcs),
each pointer carries an extra bit indicating negation.  This makes
$\neg f$ a constant-time operation and avoids storing both $f$ and
$\neg f$.

\textbf{Memory savings:}  At most $2\times$ in node count, typically
15--30\% in practice.  The bigger benefit is performance: $O(1)$
negation eliminates the recursive \texttt{Not} traversals that occur
during variable substitution (the \texttt{bdd-compose} function calls
\texttt{Not} once per substitution, and the reachability loop
performs 32 substitutions per iteration).

\textbf{Implementation cost:}  1--2 weeks.  Complement edges are
invasive to retrofit: every BDD operation must propagate complement
bits, the unique table must normalize for them, and Modula-3's
type system (where \texttt{BDD.T} is a traced reference) does not
offer the pointer-bit tricks available in C.  One would need either
wrapper objects (adding per-reference overhead) or unsafe
\texttt{LOOPHOLE} operations.

\textbf{Assessment:}  Moderate improvement for high implementation
cost.  Worth doing only as part of a broader effort to make the BDD
library production-quality.

%======================================================================
\section{Approach~6: Saturation}
%======================================================================

Saturation~\cite{ciardo2001} is a reachability algorithm designed
specifically for asynchronous systems with locality.  It replaces
BDDs with multi-valued decision diagrams (MDDs), where each MDD
level corresponds to one process instance and each edge represents
a local state of that process.

The key insight is \emph{level-by-level firing}: instead of global
BFS, the algorithm recursively saturates each MDD level by firing
all applicable local transitions until a local fixed point is reached,
before propagating to the next level.  Cross-level events (channel
synchronizations) are fired only when both participating levels are
saturated.

\textbf{Expected benefit:}  100--1000$\times$ improvement in both
time and memory for asynchronous systems, based on published
benchmarks.

\textbf{Implementation cost:}  2--4 weeks.  Requires a new MDD
data structure (the BDD library cannot be reused), Kronecker-style
transition encoding, and the recursive saturation algorithm with
its node caching and level management.

\textbf{Assessment:}  The highest-payoff approach for scaling to
large systems ($N = 20$+), but also the highest implementation cost.
Recommended only after the simpler approaches (partitioning, frontier)
have been exhausted.

%======================================================================
\section{Comparison}
%======================================================================

\begin{center}
\begin{tabular}{lcccc}
\toprule
\textbf{Approach} &
\textbf{Memory} &
\textbf{Effort} &
\textbf{BDD lib} &
\textbf{Priority} \\
& \textbf{reduction} & & \textbf{changes?} & \\
\midrule
Frontier-based       & 2--5$\times$     & 0.5--1 day  & No  & Highest \\
Partitioned TR       & 5--20$\times$    & 2--4 days   & No  & High \\
BMC with SAT         & $\gg$10$\times$  & 3--5 days   & No  & Medium \\
Static reordering    & 1.1--1.3$\times$ & 1--2 days   & No  & Low \\
Complement edges     & 1.2--1.5$\times$ & 1--2 weeks  & Yes & Low \\
Saturation (MDD)     & 100--1000$\times$& 2--4 weeks  & Yes & Future \\
Dynamic sifting      & 1.1--1.5$\times$ & 1--2 weeks  & Yes & Low \\
\bottomrule
\end{tabular}
\end{center}

The frontier-based and partitioned approaches are complementary:
frontier reduces the size of the left operand in the conjunction,
while partitioning reduces the size of the right operand and enables
early quantification.  Combined, they should reduce peak memory by
10--50$\times$.

%======================================================================
\section{Recommended Plan}
%======================================================================

\begin{description}
\item[Phase~A (1--2 days).]
  Frontier-based image computation.  Modify the reachability loop
  to track the frontier separately and compute
  $\mathit{Image}(F_k, T)$ instead of $\mathit{Image}(R_k, T)$.
  Call \texttt{BDD.CollectGarbage} between iterations to reclaim
  dead frontier and image BDDs.

\item[Phase~B (2--4 days).]
  Partitioned transition relation with early quantification.
  Keep per-process and per-channel partitions instead of monolithic~$T$.
  Rewrite the image computation to iterate over partitions.
  Combines with Phase~A for multiplicative savings.

\item[Phase~C (3--5 days, optional).]
  BMC with SAT for deadlock detection.  Complement the BDD-based
  checker with a SAT-based engine that finds deadlocks quickly and
  with negligible memory, at the cost of not being able to prove
  freedom.
\end{description}

Phases~A and~B together should bring peak memory for $N = 8$ dining
philosophers from 2.7\,GB to well under 500\,MB, and enable scaling
to larger systems that are currently infeasible.

%======================================================================
\section{References}
%======================================================================

\begin{thebibliography}{99}

\bibitem{burch1991}
J.~R. Burch, E.~M. Clarke, D.~E. Long,
``Symbolic Model Checking with Partitioned Transition Relations,''
\emph{Proc.\ VLSI}, 1991.

\bibitem{ciardo2001}
G.~Ciardo, G.~L\"uttgen, R.~Siminiceanu,
``Saturation: An Efficient Iteration Strategy for Symbolic State-Space Generation,''
\emph{Proc.\ TACAS}, LNCS~2031, pp.~328--342, 2001.

\bibitem{biere1999}
A.~Biere, A.~Cimatti, E.~M. Clarke, Y.~Zhu,
``Symbolic Model Checking without BDDs,''
\emph{Proc.\ TACAS}, LNCS~1579, pp.~193--207, 1999.

\bibitem{rudell1993}
R.~Rudell,
``Dynamic Variable Ordering for Ordered Binary Decision Diagrams,''
\emph{Proc.\ ICCAD}, pp.~42--47, 1993.

\bibitem{force}
F.~A. Aloul, I.~L. Markov, K.~A. Sakallah,
``FORCE: A Fast and Easy-To-Implement Variable-Ordering Heuristic,''
\emph{Proc.\ GLSVLSI}, pp.~116--119, 2003.

\bibitem{mcmillan1993}
K.~L. McMillan,
\emph{Symbolic Model Checking},
Kluwer Academic Publishers, 1993.

\bibitem{somenzi1999}
F.~Somenzi,
``Binary Decision Diagrams,''
in \emph{Calculational System Design}, IOS Press, 1999.

\end{thebibliography}

\end{document}
